% Targeted resume: Google, "Software Engineer".
% Framing follows Google's resume guidance: one page, quantified impact in the
% "accomplished X, measured by Y, by doing Z" form, CS fundamentals and languages
% up front, education and projects included, no photo or personal data.
\documentclass[10pt,a4paper]{article}

\usepackage[T1]{fontenc}
\usepackage[utf8]{inputenc}
\usepackage[english]{babel}
\usepackage[left=1.3cm,right=1.3cm,top=0.9cm,bottom=0.9cm]{geometry}
\usepackage{enumitem}
\usepackage{xcolor}
\definecolor{linkblue}{RGB}{0,51,120}
\usepackage[colorlinks=true,urlcolor=linkblue,linkcolor=black,citecolor=black]{hyperref}

\pagestyle{empty}
\hyphenpenalty=10000
\exhyphenpenalty=10000
\raggedright
\setlength{\parindent}{0pt}
\setlength{\parskip}{0pt}

\newcommand{\cvsection}[1]{%
  \vspace{5pt}{\large\bfseries #1}\par\vspace{2pt}\hrule\vspace{3pt}}

% #1 -- job title, #2 -- dates, #3 -- organization and location
\newcommand{\cventry}[3]{%
  \textbf{#1} \textbar{} #3 \textbar{} #2\par\vspace{1pt}}

\newlist{cvitems}{itemize}{1}
\setlist[cvitems]{label=\textbullet, leftmargin=1.2em, topsep=2pt, itemsep=0pt, parsep=0pt}

\begin{document}

{\LARGE\bfseries Kirill Kozlov}\par\vspace{3pt}
Software Engineer (Backend, Distributed Services, Applied AI)\par\vspace{3pt}
Moscow, Russia (GMT+3), open to relocation\par
Phone: +7 (910) 968 1153 | Email: \href{mailto:kirillkozlov.yar@gmail.com}{kirillkozlov.yar@gmail.com}\par
GitHub: \href{https://github.com/noobmaster-ru}{github.com/noobmaster-ru}

\cvsection{Summary}
Software engineer and Machine Learning master's student at CMC MSU with production experience at X5 Tech. Built Python services using FastAPI, PostgreSQL, and Kubernetes; developing AI product search and catalog-processing pipelines for Korzina.

\cvsection{Technical Skills}
\textbf{Programming languages:} Python, C++, Go, SQL, TypeScript\par
\textbf{Fundamentals:} data structures and algorithms, system design, REST API design, asynchronous programming,
relational data modeling, testing (pytest), code review\par
\textbf{Backend and infrastructure:} FastAPI, SQLAlchemy (async), Alembic, Pydantic, asyncio, httpx, PostgreSQL, Redis,
Apache Kafka, Docker, Kubernetes (K8s), Helm, GitLab CI/CD, Grafana, Linux, Git\par
\textbf{Applied AI:} Large Language Models (LLM), agentic workflows, RAG, prompt engineering, LangChain, OpenAI API,
OCR (Optical Character Recognition), A/B testing, PyTorch\par
\textbf{Frontend:} React, TypeScript\par
\textbf{Spoken languages:} Russian --- Native, English --- B2

\cvsection{Work Experience}
\cventry{Software Engineer (Junior)}{October 2026 --- Present}{X5 Tech (X5 Group, Russia's largest grocery retailer), Logistics Solutions Department, Moscow}
\begin{cvitems}
    \item Developing the system that handles 90\% of all foreign supplies X5 Retail Group purchases for its retail chains.
\end{cvitems}
\vspace{2pt}
\cventry{Software Engineer (Intern)}{April 2026 --- October 2026}{X5 Tech, AI Agents Department, Moscow}
\begin{cvitems}
    \item Cut legal opinion preparation time from about 97 minutes to 8--19 minutes (8 FTE saved for X5 Group) by building an agentic LLM contract-review system end to end: FastAPI and PostgreSQL backend, React and TypeScript frontend, multi-agent LangChain pipeline (document processing, requisite extraction, counterparty and property checks, risk analysis, legal opinion) with contract sections analyzed in parallel; over 500 reviews in production in May--August 2026.
    \item Kept minute-long OCR and LLM calls reliable in production by adding per-request budgets and timeouts, per-agent latency metrics in Grafana, and Kubernetes dev and prod environments deployed via Helm and GitLab CI/CD.
    \item Grounded legal analysis in Russian legislation by building a RAG source over 40 federal laws (366 articles, BM25 index) with relevant articles injected into agent prompts.
    \item Shipped an OCR and LLM service that reads foreign invoices for X5 Tech: from first commit in June 2026 to dev and prod deployment and a pilot in August 2026.
    \item Saved 10 FTE on utility-bill processing across X5 Group premises by developing the FastAPI backend that consumes OCR results from Apache Kafka, stores them and serves them to the frontend.
\end{cvitems}

\cvsection{Projects}
\cventry{\href{https://github.com/noobmaster-ru/korzina-market}{Korzina} --- Marketplace and AI Product Search}{September 2026 --- Present}{TypeScript, Node.js, PostgreSQL, React}
\begin{cvitems}
    \item Building a Telegram Mini App and bot with a catalog of 100,000+ products; automated supplier-data ingestion, normalization, and deduplication.
    \item Implemented PostgreSQL full-text and trigram search; developing a hybrid retrieval prototype combining multimodal Gemini embeddings, pgvector, reciprocal rank fusion (RRF), and Qwen reranking.
    \item Implemented LLM/VLM attribute extraction and product-compatibility checks in the prototype, with versioned embedding caches, concurrent-request deduplication, and incremental refresh.
    \item Built offline relevance evaluation with frozen datasets, qrels, Recall@20, and Precision@5; compared text and multimodal retrieval under model-cost and latency constraints.
    \item Built the Fastify API and background jobs; deployed the service using PostgreSQL, Valkey, S3, Docker, and GitLab CI/CD.
\end{cvitems}
\textbf{Other projects:} \href{https://github.com/noobmaster-ru/axiomai}{AxiomAI}: LLM-powered Telegram bot, 80\% reduction in manager workload; \href{https://github.com/noobmaster-ru/tg-mini-app-backend}{WB Parser}: Wildberries parser and Mini App, 220 users; \href{https://github.com/noobmaster-ru/wb_tables}{Placa}: analytics for 300 sellers.\par

\cvsection{Education}
\cventry{Lomonosov Moscow State University, Master's degree in Machine Learning}{September 2026 --- Present}{Faculty of Computational Mathematics and Cybernetics (CMC), Dept. of Mathematical Methods of Forecasting}
\vspace{2pt}
\cventry{Lomonosov Moscow State University, Bachelor's degree in Computer Science}{August 2022 --- August 2026}{Faculty of Computational Mathematics and Cybernetics (CMC)}
\begin{cvitems}
    \item Deep Learning School (DLS), Computer Vision and NLP tracks --- MIPT School of Deep Learning.
    \item Summer School on A/B Testing, X5 Tech --- Summer 2026 (\href{https://drive.google.com/file/d/1mUnPwkYkO6VxMJelwePo0Ie6PKA6IBot/view?usp=sharing}{certificate no. 9074503}).
\end{cvitems}

\end{document}
